% ============================================================================
% HW3 — Formal Semantics of the Ontology Vocabulary and the SHACL Layer
% TA INTERNAL — DO NOT DISTRIBUTE TO STUDENTS.
% Companion to the TA edition of the assignment handout (main.tex / TA.pdf).
% Build:  pdflatex ta-formal-semantics.tex   (twice, for cross-references)
% ============================================================================
\documentclass[11pt,a4paper]{article}

\usepackage[margin=2.6cm]{geometry}
\usepackage{amsmath,amssymb,amsthm}
\usepackage{booktabs}
\usepackage{array}
\usepackage{enumitem}
\usepackage{xcolor}
\usepackage[colorlinks=true,linkcolor=blue!50!black,urlcolor=blue!50!black]{hyperref}

\newtheorem{definition}{Definition}[section]
\newtheorem{remark}[definition]{Remark}
\newtheorem{lemma}[definition]{Lemma}
\newtheorem{convention}[definition]{Convention}

% ---- shorthand macros ------------------------------------------------------
\newcommand{\IRI}[1]{\textsf{#1}}                    % IRIs / vocabulary terms
\newcommand{\lit}[1]{\text{``}#1\text{''}}           % literals
\newcommand{\graph}{G}
\newcommand{\shapes}{\mathcal{S}}
\newcommand{\Deltares}{\Delta}
\newcommand{\dtype}{\mathrm{dt}}
\newcommand{\vio}{\mathrm{viol}}
\newcommand{\flags}{\mathrm{flags}}
\newcommand{\dbl}{\IRI{xsd{:}double}}
\newcommand{\seteq}{\stackrel{\text{set}}{=}}

\title{\textbf{HW3 Formal Semantics}\\[2pt]
\large Ontology Vocabulary and SHACL Validation, Stated Mathematically\\[4pt]
\normalsize NYCU Physical AI / TAICA --- TA INTERNAL, DO NOT DISTRIBUTE}
\author{Course Staff}
\date{Version 2026-fall-3.0 (matches \texttt{hw3-ontology.ttl} and the released shapes)}

\begin{document}
\sloppy
\maketitle

\noindent
This note restates, with mathematical precision, what the course materials define
operationally: the vocabulary of \texttt{semantic/ontology/hw3-ontology.ttl}
(\S\ref{sec:vocab}--\S\ref{sec:wf}) and the meaning of the SHACL layer
(\S\ref{sec:shacl}--\S\ref{sec:scoring}). Everything here is checkable against the
released files; nothing introduces new requirements.

% ============================================================================
\section{Preliminaries: graphs, predicates, literals}
% ============================================================================

\begin{definition}[Data graph]
An RDF \emph{data graph} $\graph$ is a finite set of triples $(s,p,o)$ over IRIs,
blank nodes, and literals. We write $\Deltares(\graph)$ for the set of subjects
and objects occurring in $\graph$ (the \emph{domain} of discourse).
\end{definition}

\begin{definition}[Predicate reading]\label{def:pred}
For a class IRI $C$ and node $x$ we write
\[
  C(x) \;\Longleftrightarrow\; (x,\ \IRI{rdf{:}type},\ C) \in \graph ,
\]
and for a property IRI $p$ and nodes/literals $x,y$,
\[
  p(x,y) \;\Longleftrightarrow\; (x,\ p,\ y) \in \graph .
\]
All quantifiers in this note range over $\Deltares(\graph)$ of the graph being
validated, and \textbf{only} over it (closed world; see \S\ref{sec:cwa}).
No inference is applied: $C(x)$ requires the \emph{explicit} typing triple, and
subclass axioms of \S\ref{sec:vocab} do \emph{not} extend the extension of $C$
during validation.
\end{definition}

\begin{definition}[Typed literals]
A literal $v$ carries a datatype $\dtype(v)$. Numeric comparison $v \le c$ is
defined for numeric datatypes (\IRI{xsd{:}double}, \IRI{xsd{:}decimal},
\IRI{xsd{:}integer}, \ldots) by their value spaces; for non-numeric $v$ the
comparison is \emph{undefined}. The course serializers emit every float as
$\lit{r}^{\wedge\wedge}\dbl$; a bare Turtle decimal such as \texttt{0.0892} has
$\dtype = \IRI{xsd{:}decimal}$ --- numerically comparable, but datatype-distinct
(this distinction matters in Def.~\ref{def:datatype}).
\end{definition}

% ============================================================================
\section{The vocabulary (signature)}\label{sec:vocab}
% ============================================================================

The signature $\Sigma = (\mathbf{C}, \mathbf{P}_o, \mathbf{P}_d)$ consists of
class names, object properties, and datatype properties. Reused terms keep their
original IRIs (declared as offline stubs); $\IRI{hw3{:}}$ mints only what no
standard covers.

\subsection{Classes}

\begin{center}\small
\begin{tabular}{@{}llll@{}}
\toprule
Class & Origin & Subclass axioms & Intended extension \\
\midrule
$\IRI{cora{:}Robot}$          & IEEE 1872 CORA & --- & the UR5 individual \\
$\IRI{soma{:}Joint}$          & SOMA & --- & kinematic joints \\
$\IRI{soma{:}RevoluteJoint}$  & SOMA & $\sqsubseteq \IRI{soma{:}Joint}$ & the six UR5 joints \\
$\IRI{soma{:}JointState}$     & SOMA & --- & one angle of one joint \\
$\IRI{soma{:}6DPose}$         & SOMA & $\sqsubseteq \IRI{dul{:}Region},\ \sqsubseteq \IRI{pos{:}Pose}$ & pose nodes \\
$\IRI{hw3{:}KinematicComputation}$ & minted & $\sqsubseteq \IRI{dul{:}Event}$ & FK/IK occurrences \\
$\IRI{hw3{:}FKComputation}$   & minted & $\sqsubseteq \IRI{hw3{:}KinematicComputation}$ & FK executions \\
$\IRI{hw3{:}IKComputation}$   & minted & $\sqsubseteq \IRI{hw3{:}KinematicComputation}$ & IK executions \\
$\IRI{hw3{:}JointConfiguration}$ & minted & --- & full-arm configurations \\
$\IRI{hw3{:}QuaternionOrientation}$ & minted & $\sqsubseteq \IRI{pos{:}Orientation}$ & $(q_x,q_y,q_z,q_w)$ nodes \\
\bottomrule
\end{tabular}
\end{center}

Subclass axioms are \emph{modelling documentation}: they align minted terms to the
upper level (DUL) and to the reused standards, but by Def.~\ref{def:pred} they play
no role in validation. The serializers therefore always assert the most specific
class expected by the shapes.

\subsection{Object properties $\mathbf{P}_o$ (with signatures $p \subseteq C \times D$)}

\begin{center}\small
\begin{tabular}{@{}lll@{}}
\toprule
Property & Signature & Super-property \\
\midrule
$\IRI{hw3{:}hasJoint}$ & $\IRI{cora{:}Robot} \times \IRI{soma{:}Joint}$ & $\sqsubseteq \IRI{dul{:}hasComponent}$ \\
$\IRI{hw3{:}computedForRobot}$ & $\IRI{hw3{:}KinematicComputation} \times \IRI{cora{:}Robot}$ & $\sqsubseteq \IRI{dul{:}involvesAgent}$ \\
$\IRI{hw3{:}solvesForTarget}$ & $\IRI{hw3{:}IKComputation} \times \IRI{soma{:}6DPose}$ & $\sqsubseteq \IRI{dul{:}hasRegion}$ \\
$\IRI{hw3{:}hasInputJointConfiguration}$ & $\IRI{hw3{:}FKComputation} \times \IRI{hw3{:}JointConfiguration}$ & \\
$\IRI{hw3{:}hasJointConfiguration}$ & $\IRI{hw3{:}IKComputation} \times \IRI{hw3{:}JointConfiguration}$ & \\
$\IRI{hw3{:}hasEndEffectorPose}$ & $\IRI{hw3{:}KinematicComputation} \times \IRI{soma{:}6DPose}$ & \\
$\IRI{hw3{:}hasJointState}$ & $\IRI{hw3{:}JointConfiguration} \times \IRI{soma{:}JointState}$ & $\sqsubseteq \IRI{dul{:}hasComponent}$ \\
$\IRI{hw3{:}isStateOfJoint}$ & $\IRI{soma{:}JointState} \times \IRI{soma{:}Joint}$ & \\
$\IRI{hw3{:}hasPositionComponent}$ & $\IRI{pos{:}Pose} \times \IRI{pos{:}Position}$ & \\
$\IRI{hw3{:}hasOrientationComponent}$ & $\IRI{pos{:}Pose} \times \IRI{pos{:}Orientation}$ & \\
$\IRI{qudt{:}hasUnit}$ & (any) $\times\ \{\IRI{unit{:}M}, \IRI{unit{:}RAD}\}$ & \\
\bottomrule
\end{tabular}
\end{center}

\subsection{Datatype properties $\mathbf{P}_d$ (with signatures $p \subseteq C \times \mathit{val}(d)$)}

\begin{center}\footnotesize
\begin{tabular}{@{}llll@{}}
\toprule
Property & Domain & Range & Read by \\
\midrule
$\IRI{hw3{:}hasDoF}$ & $\IRI{cora{:}Robot}$ & \IRI{xsd{:}integer} & \textsf{STRUCTURE} \\
$\IRI{hw3{:}jointIndex}$ & $\IRI{soma{:}Joint}$ & \IRI{xsd{:}integer} & \textsf{STRUCTURE} \\
$\IRI{hw3{:}dh\_a},\ \IRI{hw3{:}dh\_d},\ \IRI{hw3{:}dh\_alpha}$ & $\IRI{soma{:}Joint}$ & \dbl & \textsf{STRUCTURE} \\
$\IRI{hw3{:}hasJointLowerLimit},\ \IRI{hw3{:}hasJointUpperLimit}$ & $\IRI{soma{:}Joint}$ & \dbl & $\Phi_{\mathrm{JL}}$ (Def.~\ref{def:jl}) \\
$\IRI{soma{:}hasJointPosition}$ & $\IRI{soma{:}JointState}$ & \dbl & $\Phi_{\mathrm{JL}}$ \\
$\IRI{hw3{:}hasIKStatus}$ & $\IRI{hw3{:}IKComputation}$ & \IRI{xsd{:}string} & informational \\
$\IRI{hw3{:}hasResidual}$ & $\IRI{hw3{:}KinematicComputation}$ & \dbl & $\Phi_{\mathrm{NC}}$ \\
$\IRI{hw3{:}hasTargetDistance}$ & $\IRI{hw3{:}IKComputation}$ & \dbl & $\Phi_{\mathrm{AR}}$ \\
$\IRI{hw3{:}hasPoseError},\ \IRI{hw3{:}hasJacobianError}$ & $\IRI{hw3{:}FKComputation}$ & \dbl & $\Phi_{\mathrm{FK}}$ \\
$\IRI{hw3{:}hasPositionX/Y/Z}$ & $\IRI{pos{:}Position}$ & \dbl & \textsf{STRUCTURE} \\
$\IRI{hw3{:}hasQuatX/Y/Z/W}$ & $\IRI{hw3{:}QuaternionOrientation}$ & \dbl & --- \\
$\IRI{hw3{:}producedBy}$ & (any) & \IRI{xsd{:}string} & provenance \\
\bottomrule
\end{tabular}
\end{center}

% ============================================================================
\section{Well-formed grounding graphs}\label{sec:wf}
% ============================================================================

The following definitions state what a correct \texttt{data.ttl} \emph{is}; the
\textsf{STRUCTURE} shapes of \S\ref{sec:structure} are their operational check.
Fix constants: number of joints $n = 6$; course thresholds
$\theta_{\mathrm{FK}} = 0.005$, $\theta_{\mathrm{AR}} = 0.90$,
$\theta_{\mathrm{NC}} = 0.02$ (metres).

\begin{definition}[Robot specification]\label{def:robot}
A node $r$ is a \emph{well-formed robot specification} iff
$\IRI{cora{:}Robot}(r)$, $\IRI{hw3{:}producedBy}(r,\cdot)$,
$\IRI{hw3{:}hasDoF}(r, 6)$, and
\begin{gather*}
  \bigl|\{\, j : \IRI{hw3{:}hasJoint}(r,j) \,\}\bigr| = 6, \\
  \forall j\ \bigl[\IRI{hw3{:}hasJoint}(r,j) \rightarrow \IRI{soma{:}RevoluteJoint}(j)\bigr],
\end{gather*}
and every such $j$ carries $\IRI{hw3{:}jointIndex}$,
$\IRI{hw3{:}dh\_a/d/alpha}$, and both limit values, each with the exact
datatype of its range column above.
\end{definition}

\begin{definition}[Joint configuration]\label{def:jc}
A node $c$ is a \emph{well-formed joint configuration} iff
$\IRI{hw3{:}JointConfiguration}(c)$ and
\begin{gather*}
  \bigl|\{\, s : \IRI{hw3{:}hasJointState}(c,s) \,\}\bigr| = 6, \\
  \forall s\ \bigl[\IRI{hw3{:}hasJointState}(c,s) \rightarrow
     \IRI{soma{:}JointState}(s) \\
  \qquad\qquad \wedge\ \exists j\, \IRI{hw3{:}isStateOfJoint}(s,j)
     \wedge \exists v\, \IRI{soma{:}hasJointPosition}(s,v) \bigr].
\end{gather*}
\end{definition}

\begin{definition}[Structured pose]\label{def:pose}
A node $\pi$ is a \emph{well-formed pose} iff $\IRI{soma{:}6DPose}(\pi)$ and
$\exists\, \rho\ [\IRI{hw3{:}hasPositionComponent}(\pi,\rho) \wedge
\IRI{pos{:}Position}(\rho) \wedge
\IRI{hw3{:}hasPositionX/Y/Z}(\rho,\cdot)]$, with an analogous orientation
component. No pose, configuration, or numeric vector is ever encoded as a
string literal (\emph{no-JSON-strings rule}).
\end{definition}

\begin{definition}[FK / IK computation records]\label{def:comp}
$x$ is a well-formed \emph{FK record} iff $\IRI{hw3{:}FKComputation}(x)$ with
exactly one input configuration (Def.~\ref{def:jc}), exactly one end-effector
pose (Def.~\ref{def:pose}), and \dbl-typed
$\IRI{hw3{:}hasPoseError}$, $\IRI{hw3{:}hasJacobianError}$, plus provenance.
$x$ is a well-formed \emph{IK record} iff $\IRI{hw3{:}IKComputation}(x)$ with
one target link, one output configuration, a status string, and \dbl-typed
$\IRI{hw3{:}hasResidual}$, $\IRI{hw3{:}hasTargetDistance}$, plus provenance.
\end{definition}

% ============================================================================
\section{SHACL validation, formally}\label{sec:shacl}
% ============================================================================

\subsection{The validation function}

\begin{definition}[Validation]\label{def:validation}
Given a shapes graph $\shapes$ and data graph $\graph$, validation computes a
finite set of \emph{results}
\begin{multline*}
  \mathrm{Report}(\shapes, \graph) \;=\;
  \{\, (x,\ m,\ p,\ v)\ : \\
  \text{focus node } x \text{ violates a constraint with
  message } m \text{ on path } p \text{ at value } v \,\},
\end{multline*}
and $\mathrm{conforms}(\shapes,\graph) \Leftrightarrow
\mathrm{Report}(\shapes,\graph) = \emptyset$. For a shape with
$\IRI{sh{:}targetClass}\ C$, the focus nodes are $\{x : C(x)\}$ (explicit typing
only, Def.~\ref{def:pred}); for each
$\IRI{sh{:}property}\,[\,\IRI{sh{:}path}\ p\,;\ldots]$ the \emph{value nodes} of
$x$ are $\{v : p(x,v)\}$. \textbf{Each} failing value node contributes one
result --- results are per-violation, not per-shape.
\end{definition}

\subsection{Constraint components used in this course}

\begin{definition}[Value range, $\IRI{sh{:}maxInclusive}\ c$]\label{def:maxinc}
\[
  \mathrm{ok}(x) \;\Longleftrightarrow\; \forall v\, \bigl[p(x,v) \rightarrow v \le c\bigr],
  \qquad
  \vio(x,v) \;\Longleftrightarrow\; p(x,v) \wedge \neg (v \le c).
\]
Note $\neg(v \le c)$, not $v > c$: an incomparable $v$ (e.g.\ a string) also
violates. The exclusive variant $\IRI{sh{:}maxExclusive}$ replaces $v \le c$ by
$v < c$; the S3 boundary records ($v = c$ exactly) separate the two.
\end{definition}

\begin{definition}[Cardinality, $\IRI{sh{:}minCount}\ m$ / $\IRI{sh{:}maxCount}\ k$]
\[
  \mathrm{ok}(x) \;\Longleftrightarrow\;
  m \;\le\; \bigl|\{\, v : p(x,v) \,\}\bigr| \;\le\; k .
\]
\end{definition}

\begin{definition}[Datatype, $\IRI{sh{:}datatype}\ d$]\label{def:datatype}
\[
  \mathrm{ok}(x) \;\Longleftrightarrow\; \forall v\, \bigl[p(x,v) \rightarrow
  v \text{ is a literal} \wedge \dtype(v) = d\bigr].
\]
This is datatype \emph{identity}: $\dtype(\texttt{0.0892}) = \IRI{xsd{:}decimal}
\ne \dbl$, hence the classic \textsf{STRUCTURE} deduction, even though the same
value passes every numeric comparison of Def.~\ref{def:maxinc}.
\end{definition}

\begin{definition}[Class, $\IRI{sh{:}class}\ C$]
$\mathrm{ok}(x) \Leftrightarrow \forall y\,[p(x,y) \rightarrow C(y)]$, with
$C(y)$ explicit as always.
\end{definition}

\begin{definition}[SPARQL constraint, $\IRI{sh{:}sparql}$]\label{def:sparql}
A query $Q(\mathit{this})$ is evaluated once per focus node $x$ with
$\mathit{this} \mapsto x$; \textbf{every solution mapping} $\mu$ of $Q$ yields
one violation $(x, m, -, \mu)$. Thus $Q$'s body must be the \emph{negation} of
the conformance condition.
\end{definition}

\subsection{Closed world}\label{sec:cwa}

\begin{convention}[Negation as failure]
Validation treats $\graph$ as a complete description: a triple not in $\graph$
is false. Contrast OWL's open-world reading, under which absence is unknown ---
and under which numeric comparison is not even expressible. Hence the course
invariant: \emph{the ontology carries no thresholds and no cardinalities; all
checking is SHACL} (TA guide \S4.4).
\end{convention}

\begin{lemma}[Vacuous conformance]\label{lem:vacuous}
If $x$ has no $p$-value, then $x$ satisfies every value-range, datatype, and
class constraint on $p$ (the universal quantifier is vacuous), and violates a
$\IRI{sh{:}minCount}\ m$ constraint with $m \ge 1$. Presence and correctness
are therefore enforced by \emph{different} shapes, and both are needed.
\end{lemma}

% ============================================================================
\section{The course shapes as formulas}\label{sec:formulas}
% ============================================================================

\subsection{Problem shapes (student \texttt{shapes.ttl})}

\begin{definition}[FK accuracy, worked example]
\[
  \Phi_{\mathrm{FK}}:\quad
  \forall x\,\forall v\ \bigl[\IRI{hw3{:}FKComputation}(x) \wedge
  \IRI{hw3{:}hasPoseError}(x,v) \;\rightarrow\; v \le \theta_{\mathrm{FK}}\bigr]
\]
with flag \textsf{FK\_INACCURATE} on each violation.
\end{definition}

\begin{definition}[Arm out of range, TODO~1]
\begin{gather*}
  \Phi_{\mathrm{AR}}:\quad
  \forall x\,\forall v\ \bigl[\IRI{hw3{:}IKComputation}(x) \wedge
  \IRI{hw3{:}hasTargetDistance}(x,v) \;\rightarrow\; v \le \theta_{\mathrm{AR}}\bigr] \\
  \text{flag \textsf{ARM\_OUT\_OF\_RANGE}}.
\end{gather*}
\end{definition}

\begin{definition}[No convergence, TODO~2]
\begin{gather*}
  \Phi_{\mathrm{NC}}:\quad
  \forall x\,\forall v\ \bigl[\IRI{hw3{:}IKComputation}(x) \wedge
  \IRI{hw3{:}hasResidual}(x,v) \;\rightarrow\; v \le \theta_{\mathrm{NC}}\bigr] \\
  \text{flag \textsf{NO\_CONVERGENCE}}.
\end{gather*}
\end{definition}

\begin{definition}[Joint limit, TODO~3, SHACL-SPARQL]\label{def:jl}
\begin{gather*}
  \Phi_{\mathrm{JL}}:\quad
  \forall s\,\forall v\,\forall j\,\forall \ell\,\forall u\
  \Bigl[\, \IRI{soma{:}JointState}(s) \wedge \IRI{soma{:}hasJointPosition}(s,v)
  \wedge \IRI{hw3{:}isStateOfJoint}(s,j) \\
  \wedge\ \IRI{hw3{:}hasJointLowerLimit}(j,\ell) \wedge
  \IRI{hw3{:}hasJointUpperLimit}(j,u)
  \;\rightarrow\; \ell \le v \le u \,\Bigr] \\
  \text{flag \textsf{JOINT\_LIMIT\_VIOLATION}}.
\end{gather*}
SHACL Core relates a focus node only to its own values; $\Phi_{\mathrm{JL}}$
joins values across two nodes ($s$ and $j$), so it is realised as a SPARQL
constraint whose \texttt{FILTER} is the negation of the consequent:
\texttt{FILTER (?v < ?lo || ?v > ?hi)} (Def.~\ref{def:sparql}). Limits are
inclusive; a non-strict filter (\texttt{<=}, \texttt{>=}) implements
$\ell < v < u$ and falsely flags the boundary record \texttt{jl\_case\_02}
($v = u$ exactly). By closed world (Lemma~\ref{lem:vacuous}), a joint state whose
joint declares no limits yields no query solution and conforms vacuously; on
\texttt{data.ttl} Def.~\ref{def:robot} guarantees limits exist.
\end{definition}

\subsection{Structure shapes (TA \texttt{ta-shapes-full.ttl})}\label{sec:structure}

The \textsf{STRUCTURE} shapes are exactly the operational encodings of
Definitions~\ref{def:robot}--\ref{def:comp}, assembled from the components of
\S\ref{sec:shacl}: e.g.\ ``a robot has exactly six revolute joints'' is
$\IRI{sh{:}minCount}\,6 \wedge \IRI{sh{:}maxCount}\,6 \wedge \IRI{sh{:}class}\,
\IRI{soma{:}RevoluteJoint}$ on path $\IRI{hw3{:}hasJoint}$; every numeric
requirement of Def.~\ref{def:comp} is $\IRI{sh{:}minCount}\,1 \wedge
\IRI{sh{:}datatype}\,\dbl$. The TA suite additionally contains the problem
shapes $\Phi_{\mathrm{AR}}, \Phi_{\mathrm{NC}}$ (for S2) and its own copy of
$\Phi_{\mathrm{JL}}$.

% ============================================================================
\section{Flags, answer key, and scoring semantics}\label{sec:scoring}
% ============================================================================

\begin{convention}[Flag extraction]\label{conv:flag}
For a result $(x, m, p, v)$, the machine-readable flag is the prefix of $m$
before the first colon. For a record identifier $\iota$ (e.g.\
\texttt{ik\_case\_09}),
\[
  \flags(\iota) \;=\; \{\, \mathrm{prefix}(m) \ :\ (x,m,p,v) \in
  \mathrm{Report}(\shapes_{\mathrm{student}},\ \graph_{\mathrm{TA}}),\
  \iota \text{ is a substring of } x \,\}.
\]
Identifiers keep a fixed digit width so that the substring relation is a
bijection onto records.
\end{convention}

\begin{definition}[S3 --- answer-key agreement]
Let $K : \mathcal{I} \to 2^{F}$ be the answer key over record set
$\mathcal{I}$ ($|\mathcal{I}| = 30$), $F$ the four flags, and let
$\mathcal{I}^{+} = \{\iota : K(\iota) \ne \emptyset\}$ ($|\mathcal{I}^{+}|=14$),
$\mathcal{I}^{0} = \mathcal{I} \setminus \mathcal{I}^{+}$
($|\mathcal{I}^{0}|=16$). A record is \emph{correct} iff
$\flags(\iota) \seteq K(\iota)$ (set equality: no false positives, no false
negatives). Then
\[
  S_3 \;=\; 8 \cdot
  \frac{|\{\iota \in \mathcal{I}^{+} : \flags(\iota) \seteq K(\iota)\}|}{|\mathcal{I}^{+}|}
  \;+\; 2 \cdot
  \frac{|\{\iota \in \mathcal{I}^{0} : \flags(\iota) \seteq K(\iota)\}|}{|\mathcal{I}^{0}|}
  \;\in\; [0, 10].
\]
\end{definition}

\begin{definition}[S1, S2, and the Task 1 score]
With $V_{\mathrm{STRUCT}}$ the set of \textsf{STRUCTURE}-prefixed results of the
TA suite on the student's \texttt{data.ttl}:
\begin{gather*}
  S_1 = \max\bigl(0,\ 12 - 2\,|V_{\mathrm{STRUCT}}|\bigr), \\
  S_2 = 2\,[\,\textsf{AR}@\mathit{mid}\,] + 2\,[\,\textsf{AR}@\mathit{far}\,]
      + 2\,[\,\mathit{near}\ \text{clean}\,] + 2\,[\,\textsf{JL} \text{ absent}\,],
\end{gather*}
where $[\cdot]$ is the Iverson bracket, and the Task 1 score is
$S_1 + S_2 + S_3 \in [0,30]$ with pass gate (exit code $0$) at $\ge 29.9$.
Note $\mathrm{conforms}$ itself is \emph{not} the gate: a correct solution's own
data legitimately contains four violations (\emph{mid}/\emph{far} genuinely lie
beyond $\theta_{\mathrm{AR}}$).
\end{definition}

\begin{remark}[Tasks 2--4 feedback bridge]
The diagnosis of \texttt{diagnose.py} evaluates
$\mathrm{Report}(\shapes_{\mathrm{student}},\ \graph_{\mathrm{exec}})$ where
$\graph_{\mathrm{exec}}$ is built, at run time, by the student's own grounding
functions from the numeric execution records of Tasks 2--4, and prints
$\flags(\iota)$ per record. Formally it is the same validation function
(Def.~\ref{def:validation}) applied to a freshly grounded graph --- no separate
mechanism.
\end{remark}

\begin{remark}[Relation to description logic]
Each $\Phi$ of \S\ref{sec:formulas} lies in the two-variable-with-counting
fragment of FOL except $\Phi_{\mathrm{JL}}$ (five variables, a genuine join),
which is exactly why it exceeds SHACL Core and OWL alike and requires the SPARQL
escape hatch. The ontology's subclass axioms are ordinary $\mathcal{ALC}$ TBox
inclusions; they document intent but, under closed-world validation without
inference, carry no operational weight.
\end{remark}

\end{document}
